%% LaTeX2e template for FEUP's Projeto Integrador

%% 
%% FEUP, JCL, Tue Jan 30 21:57:48 2024
%%
%% PLEASE send improvements to jlopes at fe.up.pt (up till 2028 :-)

\documentclass[11pt,a4paper]{report}

%%------------------------------- preamble ------------------------------------

%% comment next for EN
\usepackage{tikz}
\usepackage[utf8]{inputenc}       % accents
\usepackage[T1]{fontenc}          % PS fonts
\usepackage{newtxtext,newtxmath}  % do not use CM fonts
\usepackage{amsmath} 
\usepackage{wrapfig}
% multi-line and other mathematical statements
\usepackage{setspace}             % setting the spacing between lines
\usepackage[normalem]{ulem}       % various types of underlining
\usepackage{caption}              % rotating captions, sideways captions, etc.
\usepackage{float}                % tables and figures in the multi-column environment 
\usepackage{subcaption}           % for subfigures and the like
\usepackage{longtable}            % tables that continue to the next page
\usepackage{multirow}             % tabular cells spanning multiple rows
\usepackage[table]{xcolor}        % driver-independent color extensions
\usepackage{lipsum}               % loren dummy text
\setlength{\marginparwidth}{2cm}  % todonotes' requirements
\usepackage{todonotes}            % todo's
\usepackage{csquotes}             % context sensitive quotation facilities
\usepackage[backend=biber,authordate]{biblatex-chicago}  % Chicago Manual of Style
\usepackage[acronym]{glossaries-extra}
\usepackage{arydshln}
\usepackage[margin=1in]{geometry}
\usepackage{xcolor}

%% document dimensions
\usepackage[a4paper,left=25mm,right=25mm,top=25mm,bottom=25mm,headheight=6mm,footskip=12mm]{geometry}
\setlength{\parindent}{1em}
\setlength{\parskip}{1ex}

%% headers & footers
\usepackage{lastpage}
\usepackage{fancyhdr}
\fancyhf{}                            % clear off all default fancyhdr headers and footers
\rhead{\small{\emph{\projtitle, \projauthor}}}
\rfoot{\small{\thepage\ de \pageref{LastPage}}} % Paginação no rodapé direito
\pagestyle{fancy}                     % apply the fancy header style
\renewcommand{\headrulewidth}{0.4pt}
\renewcommand{\footrulewidth}{0.4pt}
\pagestyle{fancy}

\usepackage{float}

%% Hiperlinks no documento
\usepackage{hyperref}
\hypersetup{
    hidelinks, % Esconde cores nos links
    linktoc=all, % Habilita links no índice
    colorlinks=false
}

%% colors
\usepackage{color}
\definecolor{engineering}{rgb}{0.549,0.176,0.098}
\definecolor{cloudwhite}{cmyk}{0,0,0,0.025}
\usepackage{booktabs} % Para tabelas mais bonitas
\usepackage{array}
%% source-code listings
\usepackage{listings}
\lstset{ %
 language=C,                        % choose the language of the code
 basicstyle=\footnotesize\ttfamily,
 keywordstyle=\bfseries,
 numbers=left,                      % where to put the line-numbers
 numberstyle=\scriptsize\texttt,    % the size of the fonts that are used for the line-numbers
 stepnumber=1,                      % the step between two line-numbers. If it's 1 each line will be numbered
 numbersep=8pt,                     % how far the line-numbers are from the code
 frame=tb,
 float=htb,
 aboveskip=8mm,
 belowskip=4mm,
 backgroundcolor=\color{cloudwhite},
 showspaces=false,                  % show spaces adding particular underscores
 showstringspaces=false,            % underline spaces within strings
 showtabs=false,                    % show tabs within strings adding particular underscores
 tabsize=2,                         % sets default tabsize to 2 spaces
 captionpos=t,                      % sets the caption-position to top
 belowcaptionskip=12pt,             % space between caption and listing
 breaklines=true,                   % sets automatic line breaking
 breakatwhitespace=false,           % sets if automatic breaks should only happen at whitespace
 escapeinside={\%*}{*)},            % if you want to add a comment within your code
 morekeywords={*,var,template,new}  % if you want to add more keywords to the set
}

%% hyperreferences (HREF, URL)
\usepackage{hyperref}
\hypersetup{
    %plainpages=false, 
    %pdfpagelayout=SinglePage,
    %bookmarksopen=false,
    %bookmarksnumbered=true,
    %breaklinks=true,
    %linktocpage,
    hidelinks,
    linktoc=all,
    colorlinks=false,
    %linkcolor=engineering,
    %urlcolor=engineering,
    %filecolor=engineering,
    %citecolor=engineering,
    %allcolors=engineering
}



%% path to the figures directory
\graphicspath{{figures/}}

%% bibliography file, must be in preamble
\addbibresource{bibliography.bib}

%% macros, to be updated as needed
\newcommand{\school}{Universidade do Minho}
\newcommand{\degree}{Licenciatura em Engenharia Informática}
\newcommand{\projtitle}{Processamento de linguagens}
\newcommand{\subtitle}{2025/2026}
\newcommand{\projauthor}{Grupo 56}
\newcommand{\projname}{Construção de um Compilador para Fortran 77 Standard}


%% my other macros, if needed
\newcommand{\windspt}{\textsf{WindsPT\/}}
\newcommand{\windscannerpt}{\emph{Windscanner.PT\/}}
\newcommand{\class}[1]{{\normalfont\sffamily #1\/}}
\newcommand{\svg}{\class{SVG}}

%% my environments for infos
\newenvironment{info}[1]{\vspace*{6mm}\color{blue}[ \textbf{INFO:} \begin{em} #1}
                        {\vspace*{3mm}\end{em} ]}
\newenvironment{infoopt}[1]{\vspace*{6mm}\color{blue}[ \textbf{INFO (elemento opcional):} \begin{em} #1}
                        {\vspace*{3mm}\end{em} ]}
                        
\renewcommand{\thesection}{\arabic{section}}


%%------------------------------- document-------------------------------------

\begin{document}


%%------------------------------- cover page ----------------------------------

\begin{titlepage}
\center

\vspace{-15mm}

\includegraphics[width=52mm]{universidade.png}

\vspace{0.5cm} 

{\large \textbf{\textsc{\school}}}\\

\vspace{0.5cm} % Espaço entre a imagem e as informações adicionais

{\large \degree}\\[8mm]

\vspace{2cm} % Espaço entre a imagem e as informações adicionais

{\Large \projtitle}\\[2mm]
{\large \subtitle}\\[28mm]

{\Large \textbf{\projauthor}}\\
\vspace{0.5cm}
{\Large \projname}\\

\vspace{0.5cm} % Espaço entre a imagem e as informações adicionais


\vspace{6cm} % Espaço entre a imagem e as informações adicionais

{\large A106804, Alice Soares} \\[2mm]
{\large A106914, Gonçalo Martins} \\[2mm]
{\large A107367, João Azevedo} \\[10mm]


{\Large Maio 2026}

\end{titlepage}

%%------------------------------- avaliação pelos pares -----------------------

\newpage


%%------------------------------- table of contents ---------------------------

%% redefine tableofcontents text, ONLY for PT
\renewcommand{\contentsname}{Índice}
\onehalfspacing
\tableofcontents
\singlespacing



\newpage

%==================================================
\section{Introdução}

O objetivo do projeto foi desenvolver um compilador para Fortran 77 standard, de acordo com o enunciado da unidade curricular de Processamento de Linguagens. O compilador deve reconhecer programas escritos em Fortran, validar a sua estrutura sintática e coerência semântica, e traduzir o código fonte para instruções executáveis na máquina virtual EWVM disponibilizada na unidade curricular.

A solução foi implementada em Python com recurso às ferramentas \texttt{ply.lex} e \texttt{ply.yacc}, seguindo uma arquitetura modular inspirada nas fases clássicas de compilação. O projeto encontra-se organizado em componentes independentes, responsáveis pelas diferentes etapas do processo de compilação, o que facilita a manutenção, os testes e a extensão futura do compilador.

Para além dos requisitos mínimos definidos no enunciado, foram também implementadas funcionalidades de valorização, incluindo criação da representação intermédia, otimizações de código, suporte a subprogramas e compatibilidade com diferentes formatos de código fonte Fortran. Estas opções não foram tratadas como acrescentos isolados, mas sim, integradas nas respetivas fases.

\section{Arquitetura e Implementação}

A arquitetura adotada e a organização do projeto foram pensadas para responder diretamente aos critérios de avaliação propostos. Segue ainda uma pipeline clássica de compilação, com separação clara entre fases. O fluxo global é:

\begin{center}
\usetikzlibrary{arrows.meta, positioning, calc}

\begin{tikzpicture}[
  node distance=0.5cm,
  box/.style={
    rectangle,
    draw,
    rounded corners=4pt,
    minimum width=2.0cm,
    minimum height=1.0cm,
    align=center,
    font=\small
  },
  arr/.style={-{Stealth[length=4pt]}, thick}
]

% Row 1
\node[box] (fortran)   {Código\\Fortran};
\node[box, right=of fortran]  (pre)  {Pré-\\processamento};
\node[box, right=of pre]      (lex)  {Análise\\léxica};
\node[box, right=of lex]      (syn)  {Análise\\sintática};
\node[box, right=of syn]      (ast)  {AST};

% Row 2 (aligned below row 1)
\node[box, below=1.2cm of fortran] (sem)  {Análise\\semântica};
\node[box, right=of sem]           (tac)  {Repr.\ inter.\\/ TAC};
\node[box, right=of tac]           (opt)  {Otimizações\\opcionais};
\node[box, right=of opt]           (ewvm) {Código\\EWVM};

% Row 1 arrows
\draw[arr] (fortran) -- (pre);
\draw[arr] (pre)     -- (lex);
\draw[arr] (lex)     -- (syn);
\draw[arr] (syn)     -- (ast);

% Bend from AST down to sem
\draw[arr] (ast.south) -- ++(0,-0.35) -- ($(sem.north)+(0,0.35)$) -- (sem.north);

% Row 2 arrows
\draw[arr] (sem)  -- (tac);
\draw[arr] (tac)  -- (opt);
\draw[arr] (opt)  -- (ewvm);

% Sublabels
\node[font=\scriptsize, below=2pt of lex]  {\texttt{ply.lex}};
\node[font=\scriptsize, below=2pt of syn]  {\texttt{ply.yacc}};

\end{tikzpicture}
\end{center}

O ponto de entrada do compilador é a classe \texttt{FortranCompiler}, responsável por coordenar as diferentes fases do processo de compilação, desde a leitura do programa Fortran até à geração de código EWVM. Esta classe também permite configurar opções relevantes, como a ativação de otimizações e a escolha entre formato fixo e formato livre. A decisão de centralizar a orquestração nesta classe evita que as fases fiquem acopladas entre si. 

A arquitetura do projeto foi organizada de forma modular, separando claramente as responsabilidades de cada etapa da compilação. Isto foi importante para manter o projeto legível e para permitir testes unitários por etapa.

NOTA: NAO SEI SE QUEREM UM DIAGRAMAZINHO OU NAO AQUI

\subsection{Pré-processamento}

A decisão de introduzir uma fase de pré-processamento responsável por normalizar o código teve dois objetivos principais.

O primeiro foi simplificar significativamente o lexer implementado com \texttt{ply.lex}, evitando regras excessivamente dependentes da posição dos caracteres na linha.

O segundo, foi permitir o suporte simultâneo aos formatos \textit{fixed-form} e \textit{free-form}. No formato fixo, remove comentários de linha, extrai labels, trata continuações e conserva o código relevante. No formato livre, remove comentários iniciados por \texttt{!} fora de strings e junta linhas continuadas com \texttt{\&}. Desta forma, o lexer torna-se praticamente independente do formato original do ficheiro, permitindo reutilizar a mesma tokenização tanto para programas Fortran 77 clássicos como para exemplos em formato moderno. O formato pode ser escolhido por linha de comandos ou configurado globalmente no compilador.

\subsection{Análise Léxica}

A análise léxica é implementada com \texttt{ply.lex}. O lexer recebe o código normalizado pelo pré-processador e converte-o numa sequência de tokens. São reconhecidas palavras-chave, identificadores, labels, literais inteiros, reais, lógicos e strings, operadores aritméticos, operadores relacionais e lógicos, delimitadores e símbolos especiais.

Outra decisão importante foi manter as palavras reservadas centralizadas na configuração do compilador. Isto facilita a extensão da linguagem reconhecida e evita duplicação entre módulos. Esta fase também produz erros estruturados quando encontra sequências inválidas, permitindo distinguir problemas de tokenização de erros sintáticos ou semânticos.


\subsection{Análise Sintática}

A análise sintática foi implementada com \texttt{ply.yacc}, de acordo com o enunciado. O parser constrói uma AST, que guarda a estrutura relevante do programa, sendo a gramática o que define as combinações válidas de tokens e permite construir essa representação estruturada do programa. Isto é essencial porque as fases seguintes precisam de conhecer a organização lógica do código, e não apenas a sequência textual em que ele foi escrito.

A análise sintática também permite rejeitar programas cuja forma não corresponde à linguagem esperada. A gramática usada pelo compilador é apresentada na Secção 4.

\subsection{Árvore Sintática Abstrata}

A AST foi desenhada com nós próprios para as construções que o compilador precisa de tratar: programa, declarações, expressões binárias e unárias, identificadores, literais, atribuições, condicionais, ciclos, I/O, arrays e subprogramas. Esta modelação evita que a análise semântica e a geração de código tenham de interpretar listas de tokens ou regras do parser, preservando apenas a informação relevante sobre o significado estrutural do programa.

Foi usado o padrão Visitor para que as operações sobre a árvore ficassem separadas da definição dos nós. Assim, a mesma AST é percorrida pelo analisador semântico e pelo gerador de IR sem que os nós conheçam os detalhes dessas fases. Adicionalmente, novas análises ou novas formas de geração de código poderiam percorrer a mesma árvore sem alterar a forma como o programa é inicialmente reconhecido. Esta decisão ajudou a manter o código modular e tornou o compilador mais extensível, facilitando a adição de novas operações sobre a árvore.

\subsection{Análise Semântica}

A análise semântica foi separada do parser para manter clara a diferença entre forma válida e significado válido. O parser aceita estruturas gramaticais e o analisador semântico decide se essas estruturas fazem sentido no contexto do programa, verificando propriedades que não podem ser garantidas apenas pela gramática. 

Nesta fase são verificadas declarações de variáveis, repetição de nomes no mesmo \textit{scope}, compatibilidade de tipos, uso de expressões lógicas em condições, coerência de labels em instruções \texttt{GOTO} e ciclos \texttt{DO}, acessos a arrays e chamadas de funções ou subrotinas. Para suportar estas validações, foi usada uma tabela de símbolos. Esta estrutura guarda informação sobre identificadores, tipos, dimensões, labels e subprogramas. A existência de scopes permite distinguir nomes definidos no programa principal de nomes definidos dentro de funções ou subrotinas.

O suporte a \texttt{FUNCTION} e \texttt{SUBROUTINE} influenciou esta fase. As assinaturas dos subprogramas são registadas antes de analisar os seus corpos, permitindo chamadas a funções definidas depois do programa principal, um padrão comum nos exemplos Fortran. Esta decisão tornou a semântica mais compatível com o estilo esperado da linguagem.

Esta etapa é fundamental para a correção do compilador, pois impede a geração de código para programas que, embora tenham uma forma sintática válida, não têm significado coerente segundo as regras da linguagem.


\subsection{Representação Intermédia}

Depois da análise semântica, a AST é traduzida para uma representação intermédia. Esta representação foi baseada em código de três endereços, uma forma linear e simplificada de representar operações do programa.

A principal razão para introduzir esta etapa foi evitar gerar código da máquina virtual diretamente a partir da AST. Embora essa abordagem fosse suficiente para uma solução mínima, tornaria mais difícil aplicar otimizações e depurar a tradução. Com uma representação intermédia, o programa passa por uma fase intermédia mais simples de analisar, transformar e converter para código final.

O código de três endereços decompõe expressões complexas em operações elementares, usando temporários para guardar resultados intermédios. Isto facilita a tradução de expressões, ciclos, condições, saltos e chamadas, que passam a ser representados por labels e instruções lineares. Também torna a geração de código EWVM menos dependente da estrutura original da árvore sintática.

Esta decisão de existência de TAC cria ainda uma base adequada para a etapa de otimização.

\subsection{Otimização}

A fase de otimização é opcional e atua sobre a representação intermédia. Foi colocada depois da geração de TAC porque nessa altura o programa já está numa forma uniforme, independente da sintaxe original, mas ainda não está preso às instruções concretas da EWVM, permitindo trabalhar sobre uma forma mais simples do programa.

As otimizações implementadas procuram reduzir redundâncias e melhorar a eficiência do código gerado, sem alterar o comportamento observável do programa. O compilador implementa os seguintes passes de otimização:

\begin{itemize}
	\item \textbf{Avaliação de constantes}: Detecta e avalia expressões cujos operandos são conhecidos em tempo de compilação, substituindo-as pelos seus resultados.
	
	\item \textbf{Propagação de constantes}: Quando uma variável recebe um valor constante, segue-se essa variável em operações subsequentes, substituindo referências pela constante quando seguro.
	
	\item \textbf{Simplificação algébrica}: Remove operações redundantes, como multiplicação por um, adição de zero ou expressões de identidade lógica.
	
	\item \textbf{Eliminação de subexpressões comuns}: Detecta sequências de instruções que produzem o mesmo resultado e guarda o resultado numa variável temporária para reutilização.
	
	\item \textbf{Remoção de código morto}: Identifica instruções que não afetam nenhuma saída observável do programa, nomeadamente atribuições a variáveis que nunca são lidas depois.
	
	\item \textbf{Simplificação de saltos}: Remove saltos incondicionais para labels imediatamente seguintes e otimiza sequências de saltos condicionais encadeados.
\end{itemize}

A decisão foi manter a otimização conservadora. Instruções com efeitos visíveis, como leituras, escritas, chamadas de subprogramas e saltos, não são removidas ou alteradas de forma agressiva. Esta escolha privilegia a correção do programa, que é mais importante do que aplicar transformações arriscadas sem uma análise de fluxo de dados mais profunda, garantindo que a melhoria de eficiência não compromete o resultado final.

\subsection{Geração de Código EWVM}

A última fase traduz a representação intermédia para código executável da EWVM, ou seja, as operações abstratas da representação intermédia são convertidas em instruções concretas da máquina virtual.

A geração de código tem de lidar com diferentes aspetos do programa, como armazenamento de variáveis, cálculo de expressões, leitura e escrita de valores, controlo de fluxo, acesso a arrays e chamadas de subprogramas. A existência da representação intermédia simplifica esta fase, porque o código já se encontra numa forma linear e próxima de instruções elementares. É mantido um mapeamento entre nomes simbólicos e posições de memória global, e cada operação TAC é convertida para uma sequência de instruções da máquina virtual.

O resultado final é um ficheiro com código VM, que pode ser executado ou inspecionado para validação manual da tradução. 

\section{Gramática}

O compilador reconhece programas Fortran 77 conformes à norma ANSI, com as limitações documentadas. A gramática foi definida tendo em conta as construções essenciais da linguagem e foi implementada através de regras em \texttt{ply.yacc}. 

O compilador suporta as seguintes categorias principais:

\begin{itemize}
	\item \textbf{Estrutura de programa}: Programa principal com opção de subprogramas (funções e subrotinas), declarações de tipo e lista de instruções.
	
	\item \textbf{Tipos de dados}: \texttt{INTEGER}, \texttt{REAL}, \texttt{LOGICAL}, \texttt{CHARACTER} e \texttt{COMPLEX}.
	
	\item \textbf{Instruções simples}: Atribuição, leitura (\texttt{READ}), escrita (\texttt{PRINT}, \texttt{WRITE}), chamada de subrotina (\texttt{CALL}), e controlo (\texttt{CONTINUE}, \texttt{GOTO}, \texttt{RETURN}, \texttt{STOP}).
	
	\item \textbf{Estruturas de controlo}: Condicional \texttt{IF/THEN/ELSEIF/ELSE/ENDIF} e ciclo \texttt{DO/ENDDO} com suporte a labels e \texttt{CONTINUE}.
	
	\item \textbf{Expressões}: Operadores aritméticos (\texttt{+}, \texttt{-}, \texttt{*}, \texttt{/}, \texttt{**}), relacionais (\texttt{.LT.}, \texttt{.LE.}, \texttt{.GT.}, \texttt{.GE.}, \texttt{.EQ.}, \texttt{.NE.}) e lógicos (\texttt{.AND.}, \texttt{.OR.}, \texttt{.NOT.}) com preced\u00eancias correctas.
	
	\item \textbf{Arrays}: Declaração e acesso a arrays unidimensionais e multidimensionais, com validação semântica de índices.
	
	\item \textbf{Funções integradas}: Suporte a funções intrínsecas como \texttt{MOD}, \texttt{ABS}, \texttt{SQRT}, \texttt{SIN}, \texttt{COS}, \texttt{EXP}, \texttt{LOG}, \texttt{INT}, \texttt{REAL}, \texttt{NINT}, \texttt{MAX} e \texttt{MIN}.
	
	\item \textbf{Subprogramas}: Definição e chamada de funções e subrotinas com parâmetros e valores de retorno.
\end{itemize}

A gramática completa em formato EBNF, incluindo todas as regras de produção, precedência de operadores, exemplos e restrições, encontra-se documentada em \texttt{docs/grammar.md}. Este documento técnico inclui também a especificação de todos os terminais (tokens), palavras-chave reservadas e limitações conhecidas do compilador. Para referência durante defesa ou estudos posteriores, recomenda-se consultar esse ficheiro.

\section{Funcionalidades Suportadas}

\subsection{Funcionalidades Essenciais}

O compilador implementa todas as construções fundamentais solicitadas no enunciado:

\begin{itemize}
	\item \textbf{Declarações de tipos}: Suporta tipos \texttt{INTEGER}, \texttt{REAL}, \texttt{LOGICAL} e \texttt{CHARACTER}, com declaração simples ou com dimensões para arrays.
	
	\item \textbf{Expressões aritméticas}: Avalia expressões envolvendo operadores \texttt{+}, \texttt{-}, \texttt{*}, \texttt{/} e \texttt{**} (exponenciação), com precedência e associatividade corretas.
	
	\item \textbf{Expressões lógicas}: Suporta operadores \texttt{.AND.}, \texttt{.OR.}, \texttt{.NOT.}, bem como avaliação de expressões booleanas em contextos apropriados.
	
	\item \textbf{Expressões relacionais}: Reconhece operadores de comparação \texttt{.EQ.}, \texttt{.NE.}, \texttt{.LT.}, \texttt{.LE.}, \texttt{.GT.}, \texttt{.GE.} para estabelecer condições.
	
	\item \textbf{Atribuições}: Suporta atribuição de valores a variáveis simples e a elementos de arrays.
	
	\item \textbf{Controlo de fluxo}: Implementa instruções \texttt{IF/THEN/ELSE/ENDIF}, ciclos \texttt{DO/CONTINUE}, \texttt{GOTO} com labels e \texttt{RETURN}.
	
	\item \textbf{I/O}: Reconhece \texttt{READ}, \texttt{PRINT} e \texttt{WRITE} com suporte a listas de variáveis e I/O dirigido.
\end{itemize}

\subsection{Funcionalidades Avançadas}

\subsubsection{Arrays}

O compilador valida se um identificador é efetivamente um array e se o número de índices é coerente na análise semântica. Na geração de código, os acessos a arrays são traduzidos através de cálculo de endereços. Para um array unidimensional, o cálculo é direto; para arrays multidimensionais, utiliza-se a fórmula de linealização em ordem linha-maior (row-major).

A geração de código está mais consolidada para arrays unidimensionais, que cobrem a maioria dos exemplos práticos fornecidos. Arrays multidimensionais são reconhecidos pela gramática mas enfrentam limitações na codificação, permanecendo como área de potencial extensão futura.

\subsubsection{Subprogramas}

O suporte a subprogramas foi integrado em todas as fases do compilador:

\begin{itemize}
	\item \textbf{Parser}: Reconhece declarações de \texttt{FUNCTION} e \texttt{SUBROUTINE}, bem como chamadas através de \texttt{CALL} ou invocação direta de funções.
	
	\item \textbf{Semântica}: Cria scopes independentes para cada subprograma, regista assinaturas (nome, tipo de retorno, parâmetros) e valida chamadas.
	
	\item \textbf{Representação Intermédia}: Representa chamadas como instruções especiais, preservando argumentos e resultados.
	
	\item \textbf{Geração de Código}: Emite instruções \texttt{CALL} ou \texttt{INVOKE} na EWVM, com passagem de argumentos e tratamento de valor de retorno.
\end{itemize}

Uma decisão importante foi registar assinaturas de subprogramas antes de analisar o corpo principal, permitindo referências a funções definidas após o programa principal, um padrão frequente em código Fortran.

\subsubsection{Otimizações}

As otimizações descritas na Secção 2.7 são aplicadas de forma opcional, ativadas por linha de comandos ou configuração. Reduzem redundâncias e melhoram a eficiência sem alterar o comportamento observável.


\section{Testes e Validação}

A validação foi organizada por fase para que uma falha indique com clareza onde está o problema. Existem testes para o lexer, parser, análise semântica, geração de IR, geração de VM, otimizador e testes end-to-end sobre programas Fortran.

Os exemplos de teste incluem os programas padrão do enunciado e alguns casos adicionais. Também existem versões em formato \textit{free-form}, usadas para validar a decisão de suportar os dois formatos de entrada.

Os testes end-to-end compilam programas reais, verificam se o código VM contém instruções essenciais como \texttt{start} e \texttt{stop}, e exportam ficheiros \texttt{.vm} para validação manual.

Para além de correr os testes com \texttt{pytest}, o projeto suporta execução com cobertura através de \texttt{pytest-cov}. Isto permite medir que partes do código são exercitadas pelos testes e ajuda a justificar a robustez da solução durante a defesa.

\section{Documentação Técnica}

Além do relatório, o projeto inclui documentação de apoio, incluindo ficheiros sobre arquitetura, gramática, decisões de design, estrutura do projeto, documentação básica de Fortran e documentação da EWVM.

O código inclui docstrings Python nas classes e funções principais, permitindo documentar responsabilidades, parâmetros, valores devolvidos e erros esperados. Esta opção torna o código mais fácil de compreender e permite gerar a documentação técnica com ferramentas Python.

A documentação em \texttt{README.md} explica de forma detalhada como instalar dependências, correr o compilador, ativar otimizações, escolher o formato de entrada, executar testes e gerar documentação, funcionando como guia rápido para utilizar o compilador.

\section{Compilação e Execução}

\subsection{Preparação do Ambiente}

Antes de compilar, é necessário preparar o ambiente:

\begin{enumerate}
	\item Executar \texttt{make setup} para criar o ambiente virtual Python e instalar dependências.
	\item Ativar o ambiente virtual com \texttt{source venv/bin/activate} (em macOS/Linux) ou \texttt{venv\textbackslash{}Scripts\textbackslash{}activate} (em Windows).
\end{enumerate}

\subsection{Compilação de Programas}

O ponto de entrada é o módulo \texttt{src.main}, invocado como:

\begin{center}
\texttt{python -m src.main <ficheiro\_fortran> [opções]}
\end{center}

As opções suportadas são:

\begin{itemize}
	\item \texttt{--format fixed|free}: Define o formato de entrada (padrão: \texttt{fixed}). Formato fixo segue Fortran 77 clássico; formato livre permite comentários com \texttt{!} e continuações com \texttt{\&}.
	
	\item \texttt{--optimize}: Ativa o pipeline de otimizações sobre a representação intermédia.
	
	\item \texttt{--output <ficheiro>}: Especifica o caminho do ficheiro de saída em código VM (padrão: \texttt{<nome\_entrada>.vm}).
\end{itemize}

\subsection{Exemplos de Uso}

\begin{itemize}
	\item Compilar um programa em formato fixo:\\
	\texttt{python -m src.main tests/examples/exemplo1\_hello.f}
	
	\item Compilar com otimizações ativadas:\\
	\texttt{python -m src.main tests/examples/exemplo1\_hello.f --optimize}
	
	\item Compilar um programa em formato livre:\\
	\texttt{python -m src.main tests/examples/free/exemplo1\_hello.f90 --format free}
	
	\item Compilar e guardar em ficheiro específico:\\
	\texttt{python -m src.main programa.f --output resultado.vm --optimize}
\end{itemize}

\subsection{Configuração Global}

É possível alterar o formato de entrada por omissão em \texttt{src/config.py}, modificando a opção global de \texttt{fixed} para \texttt{free}. Esta configuração é útil quando se pretende compilar vários ficheiros do mesmo formato sem repetir a opção na linha de comandos.

\subsection{Automatização com Makefile}

O projeto fornece um \texttt{Makefile} com atalhos para operações comuns:

\begin{itemize}
	\item \texttt{make setup}: Prepara o ambiente (cria \texttt{venv}, instala dependências).
	
	\item \texttt{make compile}: Compila um exemplo em formato fixo.
	
	\item \texttt{make compile-opt}: Compila com otimizações.
	
	\item \texttt{make compile-free}: Compila um exemplo em formato livre.
	
	\item \texttt{make compile-free-opt}: Compila em formato livre com otimizações.
	
	\item \texttt{make test-all}: Executa todos os testes.
	
	\item \texttt{make test-<fase>}: Executa testes de uma fase específica (\texttt{lexer}, \texttt{parser}, \texttt{semantic}, \texttt{ir\_codegen}, \texttt{vm\_codegen}, \texttt{optimizer}, \texttt{compiler}).
\end{itemize}


\section{Dificuldades Encontradas}

Uma das dificuldades principais foi o formato fixo de Fortran 77. Como o significado de uma linha depende da posição dos caracteres, implementar tudo diretamente no lexer tornaria as regras de \texttt{ply.lex} demasiado complexas. A solução foi criar um pré-processador responsável por labels, comentários e linhas de continuação. Esta decisão simplificou o lexer e permitiu suportar também o formato livre.


Outra dificuldade foi a distinção entre chamada de função e acesso a array. Em Fortran, uma construção como \texttt{F(X)} pode representar uma chamada à função \texttt{F} ou um acesso ao array \texttt{F}. Esta ambiguidade não pode ser resolvida apenas pela sintaxe. A solução passou por construir uma representação inicial e usar informação semântica sobre declarações e subprogramas para interpretar corretamente cada caso.


A validação de labels também exigiu cuidado. O enunciado refere explicitamente a coerência entre labels dos ciclos \texttt{DO} e comandos \texttt{CONTINUE}. Para isso, o analisador semântico regista labels definidos, labels usados por \texttt{GOTO} e labels associados a ciclos, detetando referências inexistentes ou inconsistências.


A geração de código para arrays e subprogramas foi outra fonte de complexidade. Arrays exigem cálculo de endereços e distinção entre o valor da variável e a posição de memória onde ela está guardada. Subprogramas exigem \textit{scopes}, parâmetros, chamadas e retorno, o que afeta tanto a semântica como a IR e a geração de VM.

Finalmente, a otimização teve de ser implementada de forma conservadora. Certas instruções não podem ser removidas ou reordenadas sem uma análise de efeitos mais completa. Por isso, os passes implementados privilegiam segurança e previsibilidade em vez de transformações agressivas.

\section{Conclusão}

O compilador desenvolvido cumpre com sucesso as fases principais propostas no enunciado: análise léxica com \texttt{ply.lex}, análise sintática com \texttt{ply.yacc}, análise semântica com tabela de símbolos, tradução para representação intermédia baseada em código de três endereços e geração de código executável para a EWVM. A arquitetura modular foi fundamental para separar responsabilidades, facilitar a manutenção e permitir testes unitários por etapa.

\subsection{Contribuições Além do Mínimo Exigido}

O projeto estende-se significativamente para além dos requisitos base:

\begin{itemize}
	\item \textbf{Suporte a múltiplos formatos}: O pré-processador permite compilar tanto código Fortran 77 em formato fixo (posição-dependente) como em formato livre moderno, alargando a compatibilidade.
	
	\item \textbf{Representação intermédia}: A introdução de TAC como camada intermédia facilita análises, transformações e depuração, afastando a geração final da estrutura sintática original.
	
	\item \textbf{Pipeline de otimizações}: Implementam-se múltiplos passes de otimização (constantes, propagação, simplificação, código morto), sem comprometer a correção.
	
	\item \textbf{Suporte a subprogramas}: Funções e subrotinas são plenamente integradas, com scopes, assinaturas e chamadas, permitindo modularidade nos programas compilados.
	
	\item \textbf{Documentação técnica}: Além do relatório, o projeto inclui documentação modular sobre arquitetura, gramática, decisões, estrutura e especificações técnicas.
	
	\item \textbf{Cobertura de testes}: Testes automatizados cobrem cada fase, com suporte para medição de cobertura via \texttt{pytest-cov}.
	
	\item \texttt{Docstrings em Python}: Todas as classes e funções principais documentam responsabilidades, parâmetros e comportamento esperado.
\end{itemize}

\subsection{Trabalho Futuro}

O compilador fornece uma base sólida para futuras extensões:

\begin{itemize}
	\item \textbf{Arrays multidimensionais}: Expandir a geração de código para arrays com múltiplas dimensões, incluindo ajustes de linealização e acesso.
	
	\item \textbf{Funções intrínsecas}: Adicionar suporte a funções built-in de Fortran (\texttt{SIN}, \texttt{COS}, \texttt{SQRT}, \texttt{ABS}, etc.).
	
	\item \textbf{Formatação de I/O avançada}: Estender \texttt{READ} e \texttt{WRITE} com especificadores de formato (\texttt{FORMAT} statements).
	
	\item \textbf{Análise de aliasing}: Implementar passes de análise mais sofisticados para otimizações mais agressivas e seguras.
	
	\item \textbf{Geração de relatórios de compilação}: Produzir ficheiros de saída detalhados (AST, TAC, estatísticas de otimização) para fins educacionais e de depuração.
\end{itemize}

\subsection{Reflexão Final}

O desenvolvimento deste compilador consolidou compreensão profunda sobre as fases clássicas de compilação e sua integração num sistema coerente. As decisões de design---modularidade, uso de padrões de conceção, validação por etapa---tornaram o projeto extensível e defensável. O suporte tanto a formato fixo como livre demonstra flexibilidade da arquitetura, enquanto o pipeline de otimizações e a representação intermédia evidenciam ambição para além do mínimo. Este compilador constitui um artefato educacional de qualidade, adequado para demonstração prática de conceitos de processamento de linguagens.


\begin{thebibliography}{9}

\bibitem{ewvm}
EWVM, \emph{Máquina Virtual usada na unidade curricular}, \url{https://ewvm.epl.di.uminho.pt/}.
\end{thebibliography}


\end{document}

